% Chapter 6: DevOps & Deployment
% BSU Engineering Portal Graduation Project

\chapter{DevOps \& Deployment}
\label{chap:devops}

The deployment philosophy for the BSU Engineering Portal rejects the traditional "works on my machine" anti-pattern. Instead, the system embraces \textbf{infrastructure as code}, \textbf{immutable artifacts}, and \textbf{automated deployment pipelines}. This chapter examines the production-grade containerization strategy using multi-stage builds, the sophisticated startup orchestration via \texttt{entrypoint.sh}, health-check integration for container monitoring, Infrastructure as Code (IaC) for multi-cloud deployment, and the CI/CD automation that ensures reliable, reproducible releases.

\section{Multi-Stage Containerization}
\label{sec:containerization}

\subsection{The Builder Pattern: Backend Dockerfile}

The production backend image is built using a two-stage Dockerfile~\cite{docker} (\texttt{backend-dockerfile}) that separates build dependencies from runtime~\cite{docker-best-practices}. This reduces the final image size by approximately 60\% and eliminates unnecessary build tools from production.

\begin{lstlisting}[float, language=Dockerfile, caption=Backend Multi-Stage Build (\texttt{backend-dockerfile})]
# Stage 1: Builder
FROM docker.io/python:3.11-slim AS builder

WORKDIR /app

# Install build dependencies
RUN apt-get update && apt-get install -y \
    pkg-config default-libmysqlclient-dev \
    build-essential && rm -rf /var/lib/apt/lists/*

COPY backend/requirements.txt .
RUN pip install --no-cache-dir --prefix=/install -r requirements.txt

# Stage 2: Runtime
FROM docker.io/python:3.11-slim

WORKDIR /app

# Install runtime dependencies only
RUN apt-get update && apt-get install -y \
    default-libmysqlclient-dev \
    && rm -rf /var/lib/apt/lists/*

COPY --from=builder /install /usr/local
COPY backend/ .

RUN chmod +x entrypoint.sh
EXPOSE 8000

CMD ["./entrypoint.sh"]
\end{lstlisting}

Key architectural decisions:
\begin{itemize}
    \item \textbf{--prefix=/install}: Isolates build artifacts in a staging directory for clean copy operations
    \item \textbf{Runtime dependency reduction}: Only \texttt{default-libmysqlclient-dev} is needed at runtime, not \texttt{build-essential}
    \item \textbf{Byte-code cleanup}: Post-install cleanup of \texttt{\_\_pycache\_} directories reduces image bloat
\end{itemize}

\subsection{Static Site Generation: Frontend Dockerfile}

The frontend container (\texttt{frontend-dockerfile}) demonstrates the canonical pattern for React~\cite{react} applications: build once with Node.js, then serve with Nginx~\cite{nginx}.

\begin{lstlisting}[float, language=Dockerfile, caption=Frontend Multi-Stage Build (\texttt{frontend-dockerfile})]
# Stage 1: Build - Compile the React app
FROM docker.io/node:20-alpine AS builder

WORKDIR /app

# Install dependencies first (leverage Docker cache)
COPY frontend/package.json frontend/package-lock.json ./
RUN npm ci --only=production=false --silent

# Copy source and build
COPY frontend/src/ ./src/
COPY frontend/public/ ./public/
COPY frontend/index.html ./
COPY frontend/vite.config.js ./

# Build for production
RUN npm run build

# Stage 2: Serve - Ultra minimal nginx image
FROM docker.io/library/nginx:alpine

# Remove default nginx config
RUN rm -rf /usr/share/nginx/html/*

# Copy only the built files
COPY --from=builder /app/dist /usr/share/nginx/html

# Copy custom nginx config
COPY nginx.conf /etc/nginx/conf.d/default.conf

RUN chown -R nginx:nginx /usr/share/nginx/html

EXPOSE 80

HEALTHCHECK --interval=30s --timeout=3s \
    CMD wget --quiet --tries=1 --spider http://localhost/ || exit 1

CMD ["nginx", "-g", "daemon off;"]
\end{lstlisting}

This design eliminates the Node.js runtime from production entirely. Only compiled static assets and Nginx (a battle-tested C server) handle user traffic, dramatically reducing the attack surface.

\subsection{Docker Buildx Bake for Cache Optimization}

The project leverages \texttt{docker-bake.hcl}~\cite{docker} for advanced build orchestration, enabling parallel multi-platform builds with intelligent caching.

\begin{lstlisting}[language=bash, caption=Docker Buildx Bake Configuration (\texttt{docker-bake.hcl})]
function "cache_to" {
    params = [image_name]
    result = CACHE_TYPE == "registry" ? 
        ["type=registry,ref=mhmdocker1/${image_name}:cache"] : 
        ["type=local,dest=${CACHE_DIR}/${image_name}"]
}

target "backend" {
    dockerfile = "backend-dockerfile"
    tags = ["mhmdocker1/bsu_backend:latest"]
    cache-to = cache_to("bsu_backend")
    cache-from = cache_from("bsu_backend")
}
\end{lstlisting}

\textbf{Cache Optimization Results:} The multi-stage build with intelligent caching reduces image rebuild times from minutes to seconds when only application code changes (not dependencies). Build cache is preserved between CI runs using GitHub Actions cache or registry-based storage.

\subsection{CI/CD Pipeline Flow Diagram}

The CI/CD pipeline flow diagram illustrates the complete pipeline from code commit to production deployment. The pipeline consists of the following stages:

* Code push: The developer pushes the code changes to the GitHub repository.
* GitHub Actions trigger: The GitHub Actions workflow is triggered by the code push event.
* Docker build: The Docker image is built using the Dockerfile.
* Trivy security scan: The Docker image is scanned for security vulnerabilities using Trivy.
* Docker Hub push: The Docker image is pushed to Docker Hub.
* Deployment: The Docker image is deployed to the production environment.

\begin{figure}[H]
    \centering
    \includegraphics[width=0.95\textwidth]{cicd_pipeline.png}
    \caption{CI/CD Pipeline Flow Diagram showing the automated build, security scan, and deployment process from code commit to production.}
    \label{fig:cicd_pipeline}
\end{figure}

\section{Production Docker Compose}
\label{sec:docker-compose-prod}

\subsection{Health-Check Driven Orchestration}

The production compose file (\texttt{docker-compose-im.yml})~\cite{docker-compose} implements service health verification before dependency startup, preventing race conditions between database initialization and application connection attempts.

\begin{lstlisting}[float, language=yaml, caption=Production Docker Compose with Healthchecks (\texttt{docker-compose-im.yml})]
services:
  db:
    image: docker.io/mysql:8.0
    container_name: bsu_db
    restart: always
    env_file:
      - .env.production
    volumes:
      - db_data:/var/lib/mysql
    healthcheck:
      test: ["CMD", "mysqladmin", "ping", "-h", "localhost"]
      interval: 30s
      timeout: 5s
      retries: 5
    networks:
      - bsu_network

  backend:
    image: docker.io/mhmdocker1/bsu_backend:stable
    container_name: bsu_backend
    ports:
      - "8001:8000"
    depends_on:
      db:
        condition: service_healthy  # Wait for DB health, not just start
    healthcheck:
      test: ["CMD", "curl", "-f", "http://localhost:8000/api/health/"]
      interval: 30s
      timeout: 10s
      retries: 3
      start_period: 40s
\end{lstlisting}

The \texttt{condition: service\_healthy} directive ensures the backend only starts after MySQL passes health checks---preventing connection failures during initial startup.

\subsection{Environment Configuration Management}

Production secrets and configuration are externalized to \texttt{.env.production}, ensuring no credentials are baked into images or committed to version control.

\begin{lstlisting}[language=bash, caption=Production Environment Configuration (\texttt{.env.production})]
# Django Settings
DEBUG=False
SECRET_KEY='k8f@vn2^jx-y4wm9$p!h3zqa(c+b7=e6rt*u_o1d&i5sg0lw%'
ALLOWED_HOSTS=localhost,127.0.0.1,backend

# Database Configuration
MYSQL_ROOT_PASSWORD=root
MYSQL_DATABASE=bsu_db
MYSQL_USER=bsu_user
MYSQL_PASSWORD=bsu_password
DB_HOST=db
DB_PORT=3306

# CORS Configuration
CORS_ALLOWED_ORIGINS=http://localhost:8081

# Redis Configuration
REDIS_URL=redis://redis:6379/1
\end{lstlisting}

\subsection{Redis Session and Cache Storage}

The production deployment uses Redis~\cite{redis} as a shared session store and caching layer, enabling horizontal scaling of multiple backend instances.

\begin{lstlisting}[float, language=yaml, caption=Redis Service Configuration]
redis:
  image: docker.io/redis:7-alpine
  container_name: bsu_redis
  restart: always
  volumes:
    - redis_data:/data
  healthcheck:
    test: ["CMD", "redis-cli", "ping"]
    interval: 10s
    timeout: 5s
    retries: 5
\end{lstlisting}

\textbf{Key Benefits:}
\begin{itemize}
    \item \textbf{Shared Session Store}: All backend instances share the same session data, enabling load balancing
    \item \textbf{Sub-millisecond Lookups}: Redis in-memory storage is 10-50x faster than database session queries
    \item \textbf{Cache Layer}: Frequently accessed data cached with configurable TTL (5-15 minutes)
    \item \textbf{Automatic Fallback}: If Redis is unavailable, Django falls back to database sessions
\end{itemize}

\section{Application Startup Orchestration}
\label{sec:startup-orchestration}

\subsection{Entrypoint Script as Init System}

The \texttt{entrypoint.sh} script functions as a container init system, executing sequential startup tasks with error handling and graceful degradation.

\begin{lstlisting}[language=bash, caption=Container Entrypoint Orchestration (\texttt{entrypoint.sh})]
#!/bin/bash
set -e

echo "Starting BSU Backend..."

# Run migrations in correct order
python manage.py migrate --noinput || { 
    echo "ERROR: migrations failed, aborting startup"; 
    exit 1; 
}

# Collect static files for production
python manage.py collectstatic --noinput || true

# Create superuser if it doesn't exist
python manage.py shell -c "
from django.contrib.auth import get_user_model
User = get_user_model()
admin, created = User.objects.get_or_create(
    username='admin',
    defaults={'is_superuser': True, 'is_staff': True, 'role': 'ADMIN'}
)
if created:
    admin.set_password(os.environ.get('ADMIN_PASSWORD', 'password123'))
    admin.save()
" || echo "Warning: admin user creation failed"

# Start gunicorn~\cite{gunicorn}
exec gunicorn bsu_backend.wsgi:application \
    --bind 0.0.0.0:${PORT:-8000} \
    --workers 2 \
    --timeout 120
\end{lstlisting}

The \texttt{set -e} flag ensures the container exits immediately if any step fails, triggering Docker's restart policy rather than running in a degraded state.

\subsection{Health Check Endpoint}

The Django application exposes a \texttt{/api/health/} endpoint that Docker uses to determine container readiness for traffic routing.

\begin{lstlisting}[language=Python, caption=Container Health Check Endpoint (\texttt{health.py})]
from django.http import JsonResponse
from django.db import connection

def health_check(request):
    """Health check for container orchestration.
    GET /api/health/ -> {"status": "ok", "db": "connected"}
    """
    try:
        with connection.cursor() as cursor:
            cursor.execute("SELECT 1")
        db_status = "connected"
    except Exception:
        db_status = "disconnected"

    status_code = 200 if db_status == "connected" else 503
    return JsonResponse(
        {"status": "ok" if db_status == "connected" else "unhealthy", 
         "db": db_status},
        status=status_code,
    )
\end{lstlisting}

\begin{figure}[H]
    \centering
    \includegraphics[width=0.7\textwidth]{sequence_healthcheck_1.png}
    \caption{Sequence Diagram showing container health check probe sequence: Docker daemon → Nginx → Django backend → MySQL database and Redis connectivity verification.}
    \label{fig:healthcheck-sequence}
\end{figure}

\section{Infrastructure as Code (IaC)}
\label{sec:iac}

The project includes Terraform~\cite{terraform} configurations for deploying to AWS and Azure, enabling infrastructure portability and version-controlled environments.

\subsection{AWS ECS Fargate Deployment}

The AWS architecture uses ECS Fargate (serverless containers) with Application Load Balancer for automatic scaling without managing EC2 instances.

\begin{lstlisting}[float, language=terraform, caption=AWS ECS Service Definition (excerpt from \texttt{ecs.tf})]
resource "aws_ecs_service" "backend" {
  name            = "bsu-engineering-portal-backend"
  cluster         = aws_ecs_cluster.main.id
  task_definition = aws_ecs_task_definition.backend.arn
  desired_count   = 1
  launch_type     = "FARGATE"

  network_configuration {
    subnets          = aws_subnet.private[*].id
    security_groups  = [aws_security_group.ecs_tasks.id]
    assign_public_ip = false
  }

  load_balancer {
    target_group_arn = aws_lb_target_group.backend.arn
    container_name   = "backend"
    container_port   = 8000
  }
}
\end{lstlisting}

Key AWS components:
\begin{itemize}
    \item \textbf{Fargate}: Serverless compute---no EC2 management required
    \item \textbf{Application Load Balancer}: HTTP routing and SSL termination
    \item \textbf{RDS MySQL}: Managed database with automated backups
    \item \textbf{VPC with private subnets}: Network isolation for database tier
\end{itemize}


\subsection{Azure Container Apps Deployment}

The Azure architecture uses Container Apps for serverless container hosting with HTTP-based auto-scaling.

\begin{lstlisting}[language=terraform, caption=Azure Container App Definition (excerpt from \texttt{container\_apps.tf})]
resource "azapi_resource" "backend" {
  type      = "Microsoft.App/containerApps@2023-05-01"
  name      = "bsu-engineering-portal-backend"
  
  body = jsonencode({
    properties = {
      environmentId = azapi_resource.environment.id
      configuration = {
        ingress = {
          external = true
          targetPort = 8000
        }
      }
      template = {
        containers = [{
          name   = "backend"
          image  = var.backend_image
          resources = {
            cpu    = 0.5
            memory = "1Gi"
          }
          env = [
            { name = "DB_HOST", value = azurerm_mysql_flexible_server.main.fqdn }
          ]
        }]
        scale = {
          minReplicas = 1
          maxReplicas = 3
        }
      }
    }
  })
}
\end{lstlisting}

Azure components provide equivalent functionality with native integration:
\begin{itemize}
    \item \textbf{Container Apps}: HTTP-scaling serverless containers
    \item \textbf{MySQL Flexible Server}: Managed database with private endpoint
    \item \textbf{Azure Container Registry}: Private image storage with managed identity
    \item \textbf{Key Vault}: Secret management for database credentials
\end{itemize}


\subsection{Cloud Cost Comparison}

\begin{table}[H]
    \centering
    \renewcommand{\arraystretch}{1.3}
    \begin{tabular}{|p{0.25\textwidth}|c|c|p{0.35\textwidth}|}
        \hline
        \textbf{Component} & \textbf{AWS} & \textbf{Azure} & \textbf{Notes} \\
        \hline
        Compute (monthly) & \textasciitilde\$20-40 & \textasciitilde\$15-30 & Fargate vs Container Apps (low traffic) \\
        \hline
        Database & \textasciitilde\$15 & \textasciitilde\$15 & db.t3.micro / B1s tier \\
        \hline
        Load Balancer & \textasciitilde\$20 & Included & ALB hourly charges \\
        \hline
        \textbf{Total (Dev)} & \textasciitilde\$60-85 & \textasciitilde\$35-55 & Azure more cost-effective for this workload \\
        \hline
    \end{tabular}
    \caption{Monthly Infrastructure Cost Comparison (Development Environment)}
    \label{tab:cloud_costs}
\end{table}

\textbf{Infrastructure Provisioning:} The Terraform configuration creates 12+ AWS resources including VPC, subnets, security groups, ECS cluster, RDS instance, and Application Load Balancer. The execution plan shows exactly what will be created before any changes are applied.

\section{CI/CD Pipeline}
\label{sec:cicd}

The GitHub Actions~\cite{github-actions} workflows implement a complete deployment pipeline with vulnerability scanning and automatic configuration updates~\cite{github-docs}.

\subsection{Vulnerability Scanning Gate}

Trivy~\cite{trivy} integration ensures no CRITICAL or HIGH severity vulnerabilities reach production.

\begin{lstlisting}[float, language=yaml, caption=Security Gate in CI Pipeline]
- name: Run Trivy vulnerability scanner
  uses: aquasecurity/trivy-action@master
  with:
    image-ref: '${{ vars.DOCKERHUB_USERNAME }}/bsu_backend:${{ steps.meta.outputs.version }}'
    format: 'table'
    exit-code: '1'  # FAIL on vulnerabilities
    vuln-type: 'os,library'
    severity: 'CRITICAL,HIGH'
\end{lstlisting}

\subsection{GitOps-Style Deployment Updates}

Upon successful image push, the pipeline automatically updates \texttt{docker-compose-im.yml} with the new image tag.

\begin{lstlisting}[float, language=yaml, caption=Automated Compose File Update]
- name: Update image tag in docker-compose
  run: |
    IMAGE="${{ vars.DOCKERHUB_USERNAME }}/bsu_backend:${{ needs.build.outputs.image_tag }}"
    yq -i ".services.backend.image = \"$IMAGE\"" ./docker-compose-im.yml

- name: Commit and push changes
  uses: stefanzweifel/git-auto-commit-action@v5
  with:
    commit_message: "chore: update backend image tag"
    file_pattern: "docker-compose-im.yml"
\end{lstlisting}

This enables declarative deployment: operators run \texttt{docker compose -f docker-compose-im.yml pull \&\& docker compose up -d} to deploy the exact images validated by CI.

\section{Reverse Proxy Architecture}
\label{sec:reverse-proxy}

\subsection{Nginx as Production Server}

The frontend container's Nginx configuration handles static file serving and API proxying.

\begin{lstlisting}[language=nginx, caption=Nginx Production Configuration (\texttt{nginx.conf})]
server {
    listen 80;
    server_name localhost;
    client_max_body_size 50M;

    # Security headers
    add_header X-Content-Type-Options "nosniff" always;
    add_header X-Frame-Options "SAMEORIGIN" always;

    root /usr/share/nginx/html;
    index index.html;

    location / {
        try_files $uri $uri/ /index.html;
    }

    location /api/ {
        resolver 127.0.0.11 valid=30s;
        set $backend backend;
        proxy_pass http://$backend:8000;
        proxy_set_header Host $host;
    }

    location /media/ {
        try_files $uri $uri/ =404;
        expires 30d;
        add_header Cache-Control "public, no-transform";
    }
}
\end{lstlisting}

The \texttt{try\_files \$uri \$uri/ /index.html} pattern enables React Router's client-side routing---any unmatched path falls back to \texttt{index.html} where React handles route resolution.

\section{Enhancements \& Recommendations}
\label{sec:devops-enhancements}

\subsection{Production Hardening Checklist}

\begin{enumerate}
    \item \textbf{HTTPS with Let's Encrypt}: Integrate certbot for automatic certificate provisioning
    \item \textbf{Centralized Logging}: Export container logs to ELK or cloud-native solutions (CloudWatch, Azure Monitor)
    \item \textbf{Database Encryption}: Enable RDS/Azure MySQL encryption at rest
    \item \textbf{Secret Rotation}: Implement automated credential rotation via cloud secret managers
\end{enumerate}

\subsection{Scaling Pathways}

\begin{itemize}
    \item \textbf{Vertical}: Increase Fargate/Container App CPU/memory allocation
    \item \textbf{Horizontal}: Configure auto-scaling rules based on CPU/memory or HTTP request count
    \item \textbf{Database}: Implement read replicas for query-heavy operations
    \item \textbf{CDN}: Add CloudFront/Azure CDN for static asset distribution
\end{itemize}

